%! Author = sriadityavedantam
%! Date = 4/8/26

\section{Riemann-Roch Theorem}

Let $X$ be a smooth projective curve.
Let $K_X$ denote the canonical divisor class.
For a divisor $D$, define
\[
    \Ell(D) = \{0\}\cup\{f\in\c(X) : \div(f) + D \geq 0\},
\]
and
\[
    \ell(D) = \dim(\Ell(D)).
\]

Recall that $\Ell(K_X) \cong \Omega[X]$.
Thus
\[
    \ell(K_X) = \dim(\Omega[X]) = g.
\]

If $X \subset \pp^2$ is a smooth curve of degree $d$, then
\[
    g = \frac{(d-1)(d-2)}{2}.
\]

If $d = 3$, then $g = 1$.
If $d = 4$, then $g = 3$.
A smooth curve of genus $2$ cannot lie in $\pp^2$, but can lie in some $\pp^n$.

\begin{theorem*}{
    Any smooth projective curve can be embedded into $\pp^3$.
}
\end{theorem*}

The projectivization satisfies
\[
    \pp(\Ell(D)) = \{D' \geq 0 : D' \sim D\}.
\]

\begin{theorem*}{
    (Riemann--Roch)
    Let $X$ be a smooth projective curve and $D$ a divisor on $X$.
    Then
    \[
        \ell(D) - \ell(K_X - D) = 1 - g + \deg(D).
    \]
}
\end{theorem*}

The difference $\ell(D) - \ell(K_X - D)$ depends only on $g$ and $\deg(D)$.
However, $\ell(D)$ itself depends on the divisor.
Let $X$ be a curve with $g > 0$.
Let $D_1 = 0$.
Then $\deg(D_1) = 0$ and $\Ell(D_1) = \c[X] = \c$.
Thus $\ell(D_1) = 1$.
Let $D_2 = p - q$ for distinct points $p,q \in X$.
Then $\deg(D_2) = 0$.

\begin{proof}
    If $\ell(D_2) > 0$, then there exists $f \in \c(X)$ with a simple zero at $p$ and a simple pole at $q$.
    Thus $f : X \to \pp^1$ has degree $1$, hence is an isomorphism.
    This contradicts $g > 0$.
    Therefore $\ell(D_2) = 0$.
\end{proof}

Thus $\deg(D_1) = \deg(D_2)$ but $\ell(D_1) \neq \ell(D_2)$.
Applying Riemann--Roch:
\[
    \ell(D_1) - \ell(K_X) = 1 - g
\]
implies
\[
    \ell(K_X) = g.
\]

\[
    \ell(D_2) - \ell(K_X - D_2) = 1 - g
\]
implies
\[
    \ell(K_X - D_2) = g - 1.
\]

Taking $D = K_X$ in Riemann--Roch:
\[
    \ell(K_X) - \ell(0) = 1 - g + \deg(K_X).
\]

Since $\ell(K_X) = g$ and $\ell(0) = 1$, we obtain
\[
    g - 1 = 1 - g + \deg(K_X).
\]

Hence
\[
    \deg(K_X) = 2g - 2.
\]

For plane curves:
\[
    \deg(K_X) = d(d - 3).
\]

Thus
\[
    d(d - 3) = 2g - 2 \implies g = \frac{(d - 1)(d - 2)}{2}.
\]

If $\deg(D) > 2g - 2$, then $\deg(K_X - D) < 0$.
Thus $\ell(K_X - D) = 0$ and
\[
    \ell(D) = 1 - g + \deg(D).
\]

\begin{theorem*}{
    Let $X$ be a smooth projective curve and $f\in\c(X)$ with $\omega\in\Omega[X]$.
    Then
    \[
        \sum_{x\in X} \operatorname{Res}_x(f\omega) = 0.
    \]
}
\end{theorem*}

Pick a local coordinate $t$ around $x$ such that
\[
    f\omega = \left(\frac{a_k}{t^k} + \cdots + \frac{a_1}{t} + a_0 + \cdots \right)dt.
\]

Then
\[
    \operatorname{Res}_x(f\omega) = a_1.
\]

\begin{proof}
    By Cauchy's formula,
    \[
        \operatorname{Res}_x(f\omega) = \frac{1}{2\pi i}\int f\omega.
    \]

    Summing over poles $p_1,\ldots,p_n$:
    \[
        \sum_{i=1}^n \operatorname{Res}_{p_i}(f\omega)
        = \frac{1}{2\pi i} \sum_{i=1}^n \int_{\gamma_i} f\omega.
    \]

    By Stokes' theorem,
    \[
        = \frac{1}{2\pi i} \int_{X \setminus \bigcup D_i} d(f\omega).
    \]

    Locally $f\omega = g(t)dt$, so
    \[
        d(g(t)dt) = g'(t)\,dt \wedge dt = 0.
    \]

    Thus the sum is $0$.
\end{proof}

\begin{theorem*}{
    If $g = 0$, then $X \cong \pp^1$.
}
\end{theorem*}

\begin{proof}
    Let $D = p$.
    Then $\deg(D) = 1 > 2g - 2 = -2$.
    Thus
    \[
        \ell(D) = 1 - 0 + 1 = 2.
    \]

    So there exists a non-constant $f \in \c(X)$ with a pole of order at most $1$.
    Thus $f$ defines an isomorphism $X \cong \pp^1$.
\end{proof}

\begin{theorem*}{
    If $g = 1$, then $K_X \sim 0$.
}
\end{theorem*}

\begin{proof}
    We have $\deg(K_X) = 2g - 2 = 0$.
    Since $\ell(K_X) = g = 1$, there exists an effective divisor linearly equivalent to $K_X$.
    The only effective divisor of degree $0$ is $0$.
    Thus $K_X \sim 0$.
\end{proof}

\begin{theorem*}{
    If $g = 1$, then $X$ is isomorphic to a cubic curve in $\pp^2$.
}
\end{theorem*}

\begin{proof}
    Fix $p \in X$.
    Then $\ell(p) = 1$, $\ell(2p) = 2$, and $\ell(3p) = 3$.

    Thus there exist functions $x \in \Ell(2p)$ and $y \in \Ell(3p)$ with poles of order $2$ and $3$ at $p$.

    The map $(x,y) : X \to \pp^2$ is an embedding.

    Consider $\Ell(6p)$.
    Its dimension is $6$.
    The functions $1, x, x^2, x^3, y, y^2, xy$ lie in $\Ell(6p)$ and are linearly dependent.

    This gives a cubic relation in $x,y$, so $X$ is a cubic curve.
\end{proof}